% ── Part IV, Chapter 3 ───────────────────────────────────────
% Lord Erasmus — Julius online, the communities he moves through,
% fixing what he finds broken
% Julius Caesar Scaliger, Pacific, late 21st century

\chapter{Lord Erasmus}

\strandmarker{Julius Caesar Scaliger}{Pacific, late 21st century}

\lettrine{O}{nline} he is Lord Erasmus.

He did not choose the name carefully, which is to
say he chose it the way he chooses most things:
by arriving at it through a process he does not
fully retrace, finding it already decided when he
looks for the decision. It suits him in a way he
does not examine. It has been his name in these
spaces for nine years and it fits the way the pod
fits, which is to say imperfectly and well.

The spaces he moves through online are not the
spaces most people move through.

He does not participate in the platforms designed
for participation, the large open forums where
the logic of visibility shapes everything that
is said and how it is said and who says it and
who responds. He found those spaces early and
left them early, not from distaste exactly but
from the recognition that they were optimised
for something he was not interested in, the way
a tool can be well-made and wrong for the work
you are trying to do.

What he is interested in is systems with problems
in them.

These exist in abundance. Every piece of software
maintained by a community has a tracker, and every
tracker is a record of what is broken and what has
been attempted and what has been abandoned and what
has been deferred and what no one has looked at in
long enough that the original context has been lost
and the problem sits there accruing comments that
circle it without resolving it.

He reads the trackers the way he reads the logs:
moving through them until something stops him.

\section{}

Tonight something has stopped him in a project
he has been following for two years without
contributing to.

The project is a mapping tool for ocean current
data, open source, maintained by a small group
of people who are not oceanographers but who
care about ocean current data for reasons of
their own that he has never fully determined and
does not require to be fully determined. The tool
works for most purposes and has a bug that has
been open for fourteen months, labelled as low
priority, accumulating comments from people who
have encountered it and described it and moved on.

The comments describe the same problem from
different angles.

He has been reading them, over two years, as
a kind of distributed evidence: each comment
a data point, each data point constraining the
space of possible explanations, the space
narrowing the same way spaces narrow when the
evidence is real and the attention is sustained.
He has not commented. Lord Erasmus does not
comment until he has something to say, which
is a policy he arrived at early and has not
revised, because commenting without something
to say consumes the attention of the people
who read the comments and produces nothing
that advances the problem.

Tonight he has something to say.

\section{}

He opens the code.

It is not the kind of code he was taught to read,
because he was not taught to read any kind of code,
but it is the kind of code he has taught himself
to read, by the usual method: finding the edge of
what he understood and pushing until it extended.
He reads it the way he reads the kelp: not through
the explicit content alone but through the behaviour,
what the code does under specific conditions, where
it does what it should and where it does something
else, the gap between intended behaviour and actual
behaviour being the location of the problem the way
the gap between a cable's actual tension and its
expected tension is the location of a problem on
the farm.

The bug is in the interaction between two functions
that were written by different people at different
times and that each work correctly in isolation and
incorrectly in combination. He can see this now,
clearly, in a way he could not have seen it two
years ago or fourteen months ago when the issue
was opened. He could not have seen it then because
he did not yet have enough of the surrounding
context to know what he was looking at.

He needed the two years.

He writes the comment slowly, the way he writes
anything that matters: not drafting and revising
but waiting until the sentence is complete before
beginning to type it, so that what appears on
the screen is already what he means rather than
an approximation that requires correction.

He explains what he found and where he found it
and why the two functions interact as they do.
He does not explain how long he has been looking
at it. He does not explain who he is or what he
does or why a manual labourer on an oceanic kelp
farm has spent two years reading the issue tracker
of a small open source mapping project.

He posts the comment as Lord Erasmus.

Then he closes the tracker and opens the logs
and returns to the western grid.

\section{}

The responses arrive while he is sleeping.

He reads them in the morning before he goes out
to the farm, still in the inner layer of the work
suit, the outer layer not yet on, the coffee doing
what coffee does.

There are four responses.

The first is from one of the project's maintainers,
a person whose handle he recognises from two years
of reading the tracker, whose comments he has come
to trust as technically precise and intellectually
honest. The response confirms his diagnosis and
asks a clarifying question about the edge case he
described.

He answers the clarifying question.

The second and third responses are from people he
does not recognise, who have also encountered the
bug and are grateful that it has been diagnosed and
say so in terms that are genuine and that he reads
without false modesty and without inflation: they
are glad the problem has been identified, as he is,
because identified problems can be fixed and
unidentified ones cannot.

The fourth response is from someone who disagrees
with his diagnosis.

He reads this one twice.

The disagreement is not correct, but it is not
stupid: it identifies a different location for
the problem, one that he considered and rejected,
and the rejection was correct but he can see,
reading the response, that he did not explain
clearly enough why he rejected it. He writes
a response to the response, more slowly than
the first comment, working through the reasoning
that he compressed too far the first time.

This is, he thinks, the part that cannot be
skipped: the person who disagrees, if they are
not stupid, is showing him where the explanation
failed, which is information about the explanation
rather than about the problem.

He takes the information.

\section{}

By the time he goes out to the farm the fix has
been proposed by the maintainer and is under review.

He will not follow this part closely. The fix
either addresses what he identified or it does
not, and if it does not the problem will reappear
in a different form and someone will open a new
issue and the process will begin again. He does
not need to be present for the resolution. He
found the problem and described it and that is
the contribution available to him from where he
stands, which is enough.

He moves through the lines.

The water is a different colour today, which means
something, which he is already reading before he
has decided to read it.

Somewhere in the code of a mapping tool for ocean
current data, two functions that did not know about
each other are about to be introduced, and the
introduction, if the fix is correct, will resolve
an interaction that has been producing wrong outputs
for as long as the functions have coexisted, which
is longer than anyone who will make the fix has
been working on the project.

He does not think about this.

He thinks: the water is this colour.

He turns toward what the colour means.

